\scanpage{89}
$X$, topos

$k$ Anneau sur $X$

$k\overset{u}{\lhook\joinrel\longrightarrow}\Phi$ mono dans $k$-Module, avec
\[
\left\{\begin{aligned}
&\text{\textcircled{a}}\quad \deleted{\unsure{\Phi}}\Psi=\operatorname{Coker}u\quad k\text{-plat}\\
&\text{\textcircled{b}}\quad\forall\, k\text{-modules }M,N,\ \text{et }i\unsure{\geq|>}0\\
&\hspace{1em}\Gamma^i(M\otimes_k\Phi)\otimes N\unsure{T|]}\ \text{est flasque}
\end{aligned}\right.
\]

\uline{Lemme} $\forall$ $k$-module $\unsure{S|B|N}$, $\exists$ \deleted{\unsure{hom|homom.}} $S\xrightarrow{u_S}\widetilde N$, qui \textcircled{a'} est un $k$-mono \textquotedblleft universel'' (reste mono par tensorisation) \deleted{\unsure{tel que}} $S$ plat $\Rightarrow\operatorname{Coker}u_{\unsure{N|S}}$ plat, et \textcircled{b'} $\forall i>0$, et $k$ modules $\deleted{N}P,M$ $\Gamma_k^i(\widetilde N\otimes_k P)\otimes_k M$ flasque

On prend $\widetilde N=\Phi\otimes_k S$, $u_S=u\otimes\mathrm{id}_S$, alors a') résulte de a), et b') de b)

$\mathscr{DR}^{**}_{pd}(\deleted{\unsure{X,}}k,u)=\Gamma^*(\Psi)\otimes_k\Lambda^*\Psi$ Anneau bigradué (degré compl + degr. ext = degré total) \inserted{alterné (pour degr. ext)}, différentiel à p.d. (sur idéal des bidegrés $>0$ \unsure{...|?}) \unsure{une|...|---} résolution de $k\{T\}$. (où $\deg\mathrm{ext}\,T=0$, $\deg\mathrm{tot}\,T=1$)
\[
\mathrm{DR}_{pd}^{**}(k,u)=\Gamma(X,\mathscr{DR}_{pd}^{**}(k,u)).
\]

Soit $\mathbb{DR}^{**}$ la cat \inserted{dérivée} des anneaux bigradués (à degr $\geq 0$) alternés diff à p.d. \deleted{\unsure{un}}\inserted{avec \deleted{\unsure{unité}}\inserted{élément} \unsure{marqué pur} bidegré $(1,0)$} On va définir un \unsure{syst.|système} transitif d'isom. dans $\mathbb{DR}^{**}$ entre les $\mathrm{DR}_{pd}^{**}(k,u)$, pour $k$ fixé et $u$ variable

Notons que si on a
\diagram{\begin{tikzcd}[ampersand replacement=\&,row sep=small]
\& \Phi\arrow[dd,"\lambda"]\\
k\arrow[ur,"u"]\arrow[dr,"u'"'] \&\\
\& \Phi'
\end{tikzcd}}
d'où $\lambda^{**}:\mathrm{DR}_{pd}^{**}(k,u)\unsure{\mathord{-}|\to}\mathrm{DR}_{pd}^{**}(k,u')$, alors $\lambda^{**}$ est un iso. dans $\mathbb{DR}^{**}$

\diagram{\begin{tikzcd}[ampersand replacement=\&,column sep=small,row sep=small]
\&\Phi'\arrow[dr,hook]\&\&\\
k\arrow[ur,hook,"u"]\arrow[dr,hook,"u'"']\&\& S\arrow[r,hook,"\text{donné par lemme}"]\arrow[d,phantom,"\Vert" description]\&N\\
\&\Phi\arrow[ur,hook]\&\Phi\amalg_k\Phi'\&
\end{tikzcd}}
d'où
\diagram{\begin{tikzcd}[ampersand replacement=\&,row sep=small]
\mathrm{DR}(u)\arrow[dr,"\sim"]\&\\
\&\mathrm{DR}(N)\\
\mathrm{DR}(u')\arrow[ur,"\sim"']\&
\end{tikzcd}}
d'où $\mathrm{DR}(u)\simeq\mathrm{DR}(u')$ dans $\mathbb{DR}^{**}$

indépendance du choix de $S\hookrightarrow N$ avec $N$ comme dans le lemme
\diagram{\begin{tikzcd}[ampersand replacement=\&,column sep=small,row sep=small]
\&\&N\arrow[dr,hook]\&\&\\
k\arrow[r,phantom,"{-}" description]\&S\arrow[ur,hook]\arrow[dr,hook]\&\&\unsure{T''|T}\arrow[r,hook,"\text{donné par lemme}"]\arrow[d,phantom,"\Vert" description]\&N''\\
\&\&N'\arrow[ur,hook]\&N\amalg_S N'\&
\end{tikzcd}}
d'où
\diagram{\begin{tikzcd}[ampersand replacement=\&,column sep=small,row sep=small]
\mathrm{DR}(u)\arrow[dr]\&\&\mathrm{DR}(N)\arrow[dr,"\sim"]\&\&\\
\&\mathrm{DR}(S)\arrow[ur]\arrow[dr]\&\&\mathrm{DR}\,N''\&\unsure{\mathord{\sim}}\\
\mathrm{DR}(u')\arrow[ur]\&\&\mathrm{DR}(N')\arrow[ur,"\sim"']\&\&
\end{tikzcd}}

Transitivité évidente par diagramme
% archive sheet number: 88
